\section*{अध्याय \ref{ch:g2:what-animals-eat} --- जंतु क्या खाते हैं}

\begin{solution}{exo:g2:what-animals-eat:1}
\omterm{def:g2:what-animals-eat:kinds}{शाकाहारी}; \omterm{def:g2:what-animals-eat:kinds}{मांसाहारी}; \omterm{def:g2:what-animals-eat:kinds}{सर्वाहारी}।
\end{solution}

\begin{solution}{exo:g2:what-animals-eat:2}
(क) \omterm{def:g2:what-animals-eat:kinds}{शाकाहारी}। (ख) \omterm{def:g2:what-animals-eat:kinds}{मांसाहारी}। (ग) \omterm{def:g2:what-animals-eat:kinds}{सर्वाहारी}। (घ) \omterm{def:g2:what-animals-eat:kinds}{शाकाहारी}। (ङ)
\omterm{def:g2:what-animals-eat:kinds}{सर्वाहारी}।
\end{solution}

\begin{solution}{exo:g2:what-animals-eat:3}
चपटे पीसने वाले --- घास जैसे \omterm{def:g1:plants-around-us:plant}{पौधे} सख़्त होते हैं और निगलने से पहले
उन्हें देर तक पीसना पड़ता है; दबोचने को कोई शिकार होता ही नहीं।
\end{solution}

\begin{solution}{exo:g2:what-animals-eat:4}
छोटी और मोटी: \omterm{def:g2:from-seed-to-plant:seed}{बीज} तोड़ना। मुड़ी और अँकुड़ीदार: माँस फाड़ना। लंबी और
पतली: छाल और दरारों से कीट निकालना।
\end{solution}

\begin{solution}{exo:g2:what-animals-eat:5}
नन्हे पौध-चीलर (माहू) --- गुबरैला उन्हें बड़ी संख्या में खाता है,
इसीलिए माली उससे प्यार करते हैं।
\end{solution}

\begin{solution}{exo:g2:what-animals-eat:6}
घास हर कौर में थोड़ा ही पोषण देती है, इसलिए गाय को पेट भरने के लिए
घंटों चरना पड़ता है; माँस भरपूर होता है, इसलिए एक अच्छा भोजन लोमड़ी को
काफ़ी देर चला देता है।
\end{solution}

\begin{solution}{exo:g2:what-animals-eat:7}
गर्मियों में केंचुए और कीट; सर्दियों में जामुन, जब ज़मीन केंचुओं के
लिए बहुत सख़्त हो जाती है। चौड़ा मेनू \omterm{def:g2:what-animals-eat:kinds}{सर्वाहारी} को एक चीज़ ख़त्म होने
पर दूसरी पर जाने देता है --- मुश्किल मौसमों में बड़ी मदद।
\end{solution}

\begin{solution}{exo:g2:what-animals-eat:8}
उत्तर खुला है। आम तौर पर: गौरैया और फ़िंच (छोटी मोटी चोंचें) \omterm{def:g2:from-seed-to-plant:seed}{बीज} लेते
हैं; चिड़ियाँ कीट और \omterm{def:g2:from-seed-to-plant:seed}{बीज} चुनती हैं; कलचिड़ी केंचुओं के लिए ज़मीन कुरेदती
है और जामुन लेती है। प्रतिज्ञप्ति का चोंच-मेनू मेल अच्छा उतरता है।
\end{solution}

\begin{solution}{exo:g2:what-animals-eat:9}
\omterm{def:g2:what-animals-eat:kinds}{मांसाहारी} का मेनू: दूसरे \omterm{def:g1:animals-around-us:animal}{जंतु}। चरण 1 और 2 --- दाँत (नुकीले दबोचने
वाले) और शिकार के औज़ार (पंजे, आगे की ओर देखती \omterm{def:g1:the-five-senses:sense}{आँखें}); चरण 4 ने उसे
नाम दिया।
\end{solution}

\begin{solution}{exo:g2:what-animals-eat:10}
लोमड़ी ख़रगोश खाती है, और ख़रगोश घास खाते हैं। घास नहीं, तो ख़रगोश
नहीं; ख़रगोश नहीं, तो भूखी लोमड़ियाँ। लोमड़ी अपने शिकार के \emph{ज़रिए}
घास पर निर्भर है --- प्रकृति में भोजन शृंखलाएँ बनाते हैं।
\end{solution}
